\chapter{Functional evaluation for vector / tensor data}
\label{chap:dev:tensor-functional}

\begin{quote}
Phase metadata. Spec dir:
\passthrough{\lstinline!specs/2026-05-05-phase-15-tensor-functional/!}.
Generalizes \passthrough{\lstinline!func::evaluate!} to
\passthrough{\lstinline!Field<Vec3d>!} input with two SFINAE-detected
callable arities (\(\mathbf u\); \(\mathbf u, \nabla\mathbf u\)) and adds a
contraction expression-template (ET) layer that fuses
\passthrough{\lstinline!Field<Mat3d>!} pointwise arithmetic into the
\passthrough{\lstinline!func::integrate!} reduction. The roadmap
acceptance bar is the closed-form match of the linear-elastic energy
\(F = \int \tfrac12\,\boldsymbol{\sigma}\!:\!\boldsymbol{\varepsilon}\,dV\)
on a periodic uniaxial dilation wave. Version target:
\passthrough{\lstinline!0.15.0!}.
\end{quote}

\subsection{1. Theory}\label{phase15-theory}

\subsubsection{Tensor-rank arithmetic on \texorpdfstring{$\mathbb{R}^3$}{R3}}\label{phase15-rank}

Phase 13 settled the rank-2 surface for tensor fields, with
\(M(i, j) = \partial_j v_i\) the velocity-gradient convention so that
\(\mathrm{tr}(\nabla \mathbf v) = \nabla\!\cdot\!\mathbf v\). Phase 15
takes the obvious next step: a functional surface that consumes a
\passthrough{\lstinline!Field<Vec3d>!} and (optionally) its rank-2
gradient. Higher derivatives of vector fields would consume rank-3 or
rank-4 outputs --- those types are deferred per the AskUserQuestion
answer.

The contractions used in this phase
(\passthrough{\lstinline!trace!},
\passthrough{\lstinline!singleContract!},
\passthrough{\lstinline!doubleContract!},
\passthrough{\lstinline!sym!}) are all rank-2-in / rank-{0,2}-out:
\(\mathrm{tr}(\mathbf M) = M_{ii}\),
\((\mathbf A \cdot \mathbf B)_{ij} = A_{ik} B_{kj}\) (matrix product),
\(\mathbf A : \mathbf B = A_{ij} B_{ij}\),
\(\mathrm{sym}(\mathbf M) = \tfrac12(\mathbf M + \mathbf M^\top)\). Einstein
summation is implicit.

\subsubsection{Isotropic linear elasticity}\label{phase15-elasticity}

The constitutive law in Lamé form is
\[
\boldsymbol{\sigma}(\boldsymbol{\varepsilon})
\;=\; \lambda\,\mathrm{tr}(\boldsymbol{\varepsilon})\,\mathbf{I}
   \;+\; 2\mu\,\boldsymbol{\varepsilon}.
\]
Strain is the symmetric part of the displacement gradient:
\(\boldsymbol{\varepsilon} = \tfrac12(\nabla\mathbf u + (\nabla\mathbf u)^\top)\).
The elastic-energy density is
\[
W(\boldsymbol{\varepsilon})
\;=\; \tfrac12\,\boldsymbol{\sigma}:\boldsymbol{\varepsilon}
\;=\; \tfrac12\!\left[\lambda\,\mathrm{tr}(\boldsymbol{\varepsilon})^2
                     + 2\mu\,\boldsymbol{\varepsilon}:\boldsymbol{\varepsilon}\right].
\]
The simplification eliminates the rank-2 stress
\(\boldsymbol{\sigma}\) from the integrand --- the right-hand side is a
linear combination of two scalar-valued contractions of
\(\boldsymbol{\varepsilon}\). The ET layer's natural unit of work is
exactly this form.

\subsection{2. Math derivation}\label{phase15-math}

\subsubsection{Closed form on a periodic uniaxial dilation wave}\label{phase15-closed}

The roadmap "stretched cubic block" geometry on a periodic box is
realized by
\(\mathbf u(x, y, z) = (\alpha\sin(k_0 x), 0, 0)\) with \(k_0 = 1\) on
\([0, 2\pi]^3\). The displacement gradient is
\[
\nabla \mathbf u
\;=\;
\begin{pmatrix}
\alpha\cos(k_0 x) & 0 & 0\\
0 & 0 & 0\\
0 & 0 & 0
\end{pmatrix},
\]
which is already symmetric, so
\(\boldsymbol{\varepsilon} = \nabla\mathbf u\),
\(\mathrm{tr}(\boldsymbol{\varepsilon}) = \alpha\cos(k_0 x)\), and
\(\boldsymbol{\varepsilon}:\boldsymbol{\varepsilon}
= \alpha^2\cos^2(k_0 x)\). Substituting into \(W\):
\[
W(x)
\;=\; \tfrac12\,\big(\lambda + 2\mu\big)\,\alpha^2\cos^2(k_0 x).
\]
Integrating over the box,
\[
F_{\mathrm{ref}}
\;=\; \tfrac12\,(\lambda + 2\mu)\,\alpha^2
\,\Big[\textstyle\int_0^{2\pi}\!\cos^2(k_0 x)\,dx\Big]
\,(2\pi)^2
\;=\; 2\pi^3\,(\lambda + 2\mu)\,\alpha^2.
\]
The result is independent of which axis carries the wave (by
symmetry); the helper
\passthrough{\lstinline!analyticalLinearElasticEnergy!} retains
\passthrough{\lstinline!axis!} as a parameter only for documenting
intent at the call site.

For \(E = 1\), \(\nu = 0.3\): \(\lambda = 0.576923\), \(\mu = 0.384615\),
\(\lambda + 2\mu = 1.346154\); at \(\alpha = 0.01\),
\(F_{\mathrm{ref}} = 2\pi^3 \cdot 1.346154 \cdot 10^{-4}
\approx 8.348\times 10^{-3}\).

\subsubsection{Convergence order}\label{phase15-convergence}

The discrete energy inherits the order of the underlying gradient
operator \(\nabla\mathbf u\). Phase 7's
\passthrough{\lstinline!diff::gradient(Field<Vec3d>)!} ships at
default order 2 (the order-4 stencil is available with explicit
\passthrough{\lstinline!<Order=4>!}; the new
\passthrough{\lstinline!func::evaluate!} overload forwards to the
default). The discrete energy therefore converges as
\(O(h^2)\) on the periodic uniaxial wave. The acceptance test gate at
\(N = 64\) asks for relative error \(< 5\times 10^{-3}\); the
convergence sweep on \(N \in \{16, 32, 64\}\) gates the rate to
\(\log_2(e_N / e_{2N}) \in [1.6, 2.4]\).

\subsection{3. Algorithm}\label{phase15-algo}

\subsubsection{File layout}\label{phase15-files}

\begin{itemize}
\tightlist
\item
  \passthrough{\lstinline!include/gridcalc/tensor/Expressions.hpp!} ---
  the contraction-ET nodes, operator overloads, factory helpers, and
  \passthrough{\lstinline!materialize!} / \passthrough{\lstinline!is\_expr!} traits.
\item
  \passthrough{\lstinline!include/gridcalc/func/Functional.hpp!} ---
  extended with the
  \passthrough{\lstinline!evaluate(Field<Vec3d>, F\&\&, Tag)!} overload.
\item
  \passthrough{\lstinline!include/gridcalc/func/Integrate.hpp!} ---
  extended with the fused
  \passthrough{\lstinline!integrate(E expr, Tag)!} overload (gated on
  \passthrough{\lstinline!tensor::expr::is\_expr\_v<E> \&\& value\_type<E> == double!}).
\item
  \passthrough{\lstinline!examples/common/elastic\_energy.hpp!} ---
  shared helpers (Lamé constants, manufactured displacement, energy
  density, closed-form reference).
\item
  \passthrough{\lstinline!examples/elastic\_energy.cpp!} --- runnable
  CLI demo with VTK output.
\item
  \passthrough{\lstinline!test/tensor\_expressions\_test.cpp!},
  \passthrough{\lstinline!test/func\_evaluate\_vector\_test.cpp!},
  \passthrough{\lstinline!test/elastic\_energy\_test.cpp!},
  \passthrough{\lstinline!test/elastic\_energy\_demo\_test.cpp!} ---
  the four new test files.
\end{itemize}

\subsubsection{The \texttt{evalAt} recursion}\label{phase15-evalat}

Each ET node exposes
\(\texttt{value\_type evalAt(int i, int j, int k) const noexcept}\). The
node holds its operands by value (composing nodes own their children)
and forwards \texttt{evalAt} down the tree:

\begin{itemize}
\tightlist
\item
  \passthrough{\lstinline!Leaf<T>!} holds a non-owning
  \passthrough{\lstinline!const Field<T>*!}; \passthrough{\lstinline!evalAt!}
  reads the wrapped field. The lifetime contract is documented in the
  Doxygen.
\item
  \passthrough{\lstinline!IdentityField!} returns
  \passthrough{\lstinline!Mat3d::Identity()!} regardless of \((i, j, k)\);
  no per-cell storage.
\item
  \passthrough{\lstinline!Scaled<E>!} stores the scalar and the inner
  expression by value; \passthrough{\lstinline!evalAt!} returns
  \(c \cdot \mathrm{inner.evalAt}(i, j, k)\).
\item
  \passthrough{\lstinline!Plus<L, R>!}, \passthrough{\lstinline!Minus<L, R>!},
  \passthrough{\lstinline!Mul<L, R>!} (scalar-only) recurse into both children.
\item
  \passthrough{\lstinline!Trace<E>!} returns
  \(\mathrm{inner.evalAt}(i,j,k).\texttt{trace()}\).
\item
  \passthrough{\lstinline!Sym<E>!} returns
  \(\tfrac12(M + M^\top)\) at the cell.
\item
  \passthrough{\lstinline!SingleContract<L, R>!} returns
  \(\mathrm{lhs} \cdot \mathrm{rhs}\) (Eigen
  \passthrough{\lstinline!Matrix3d::operator*!}).
\item
  \passthrough{\lstinline!DoubleContract<L, R>!} returns
  \((\mathrm{lhs}.array() * \mathrm{rhs}.array()).sum()\).
\end{itemize}

A static-assert on each combinator gates the
\passthrough{\lstinline!value\_type!} of its operands so a
mismatched-rank composition (e.g.\ adding a
\passthrough{\lstinline!double!}-valued and a
\passthrough{\lstinline!Mat3d!}-valued expression) fails at the call
site with a readable message rather than a depth-N template
instantiation trace.

\subsubsection{The fused \texttt{func::integrate(expr, tag)}}\label{phase15-fused}

Implementation:

\begin{verbatim}
template <class E, class Tag = Pairwise,
          std::enable_if_t<tensor::expr::is_expr_v<E>
                            && std::is_same_v<typename std::decay_t<E>::value_type, double>,
                           int> = 0>
inline double integrate(const E& expr, Tag tag = {}) {
    return integrate(tensor::expr::materialize(expr), tag);
}
\end{verbatim}

The forwarding-to-scalar-overload is the simplest correct
implementation: it allocates one
\passthrough{\lstinline!Field<double>!} for the scalar integrand, then
hands that to the existing pairwise / Neumaier reducer. The
\emph{point} of the overload is what it does \emph{not} allocate ---
no \passthrough{\lstinline!Field<Mat3d>!} is built for
\(\boldsymbol{\sigma}\), \(\boldsymbol{\varepsilon}\), or any
intermediate combination thereof. That's a 9\(\times\) memory savings
per intermediate compared with the pre-Phase-15 eager path.

\subsubsection{Algorithmic complexity}\label{phase15-complexity}

\(O(N^3)\) per cell evaluation regardless of expression depth (the
\passthrough{\lstinline!evalAt!} recursion is constant-depth at
compile time, and Eigen's
\passthrough{\lstinline!Matrix3d::operator*!} /
\passthrough{\lstinline!.array().sum()!} are constant per cell).
\(O(N^3)\) memory for the scalar integrand allocated by the fused
\passthrough{\lstinline!integrate!} overload. The pre-ET-layer eager
path on \(\tfrac12(\lambda + 2\mu)\,\boldsymbol{\sigma}:\boldsymbol\varepsilon\)
allocates two extra \passthrough{\lstinline!Field<Mat3d>!} blocks
(one for \(\boldsymbol{\sigma}\), one for \(\boldsymbol{\varepsilon}\)) ---
each \(9\times\) the integrand's memory --- so the ET layer's working
set is \(\sim 19\times\) smaller for the elastic-energy demo.

\subsection{4. Design decisions}\label{phase15-design}

The four AskUserQuestion answers from the spec round, plus the
periodic-uniaxial-wave substitution, are recorded in
\passthrough{\lstinline!specs/2026-05-05-phase-15-tensor-functional/requirements.md!}
under "Decisions made for this feature." Brief reprise:

\begin{itemize}
\tightlist
\item
  \textbf{ET layer lands now.} The user requested the ET surface
  alongside the elastic-energy demo, even though the rank-2-only
  inputs do not strictly trigger the rank-6 explosion the ET layer was
  originally proposed to fend off (Phase 11 / Phase 13 deferral). The
  scope here is bounded: ten nodes, all rank-2-or-lower, with
  Phase 13's eager primitives kept untouched in parallel.
\item
  \textbf{Rank-3 / rank-4 tensor types deferred.}
  \(\boldsymbol\varepsilon = \mathrm{sym}(\nabla\mathbf u)\) is the
  highest rank we touch; the integrand reduces to two
  rank-2-contraction terms before any rank-3 surface enters the
  picture. The Eigen-unsupported-tensor-module portability question
  stays deferred.
\item
  \textbf{Demo material parameters: standard \(E = 1\), \(\nu = 0.3\),
  \(\alpha = 0.01\)} on the standard \([0, 2\pi]^3\) box. The
  closed-form constant \(F_{\mathrm{ref}} = 2\pi^3(\lambda + 2\mu)\alpha^2\)
  drops out; the convergence test gate sits in
  \([1.6, 2.4]\) for the \(O(h^2)\) rate.
\item
  \textbf{Demo binary ships, not test-only.} Mirrors Phase 12 / 14;
  three committed PNG snapshots (small / medium / large \(\alpha\))
  embedded in the User Guide chapter.
\item
  \textbf{Periodic-uniaxial-wave substitution.} A literal rigid
  uniaxial stretch \(u_x = \alpha\,x\) is not periodic on
  \([0, 2\pi]^3\) and breaks the only-\passthrough{\lstinline!Periodic!}
  policy. The dilation wave \(u_x = \alpha\sin(k_0 x)\) is the natural
  periodic substitute and admits the same \((\lambda + 2\mu)\) tensor
  signature as the rigid stretch in the closed-form energy. Phase 19
  unlocks the literal geometry. Documented in the User Guide chapter
  and in the test file's top-of-file comment so future contributors
  don't re-litigate.
\item
  \textbf{Scalar \(\times\) scalar in the ET layer, no
  \passthrough{\lstinline!Mat3d \(\times\) Mat3d!}.} The
  \passthrough{\lstinline!*!} operator between two
  \passthrough{\lstinline!double!}-valued expressions is unambiguous;
  between two \passthrough{\lstinline!Mat3d!}-valued expressions it
  could mean either matrix product or full contraction. The dedicated
  \passthrough{\lstinline!singleContract!} /
  \passthrough{\lstinline!doubleContract!} factories disambiguate at
  the call site.
\end{itemize}

\subsection{5. References}\label{phase15-refs}

\begin{itemize}
\tightlist
\item
  Gurtin, M. E., Fried, E., \& Anand, L. (1981). \emph{An Introduction
  to Continuum Mechanics}. Academic Press. ISBN 0-12-309750-9. Linear
  isotropic elasticity, the Lamé form, and the symmetric-strain
  derivation are in chapters 9--10.
\item
  Holzapfel, G. A. (2000). \emph{Nonlinear Solid Mechanics: A Continuum
  Approach for Engineering}. Wiley. ISBN 978-0-471-82319-3. Used for
  the chosen velocity-gradient index convention and the symmetric-stress
  consistency arguments; chapter 6 covers linear elasticity as a
  small-strain limit.
\item
  Vandevoorde, D., Josuttis, N. M., \& Gregor, D. (2017).
  \emph{C++ Templates: The Complete Guide}, 2nd ed. Addison-Wesley.
  ISBN 978-0-321-71412-1. Chapter 27 covers expression templates with
  modern C++17 idioms (deduction guides, ADL-friendly operator
  overloads, \passthrough{\lstinline!std::enable\_if\_t!} gates), the
  template machinery the ET layer here is built on.
\item
  Veldhuizen, T. (1995). "Expression templates." \emph{C++ Report},
  7(5), 26--31. \texttt{https://doi.org/10.1145/216973.216979}. The
  founding paper for the technique.
\item
  Eigen project documentation: \texttt{https://eigen.tuxfamily.org/dox/}
  --- in particular the \passthrough{\lstinline!Matrix3d!} surface
  used pointwise inside each
  \passthrough{\lstinline!evalAt!} call, and the matrix /
  array-arithmetic semantics. Eigen is itself a sophisticated
  expression-template library; the ET layer here piggybacks on Eigen's
  per-cell expression machinery rather than reimplementing it.
\item
  LeVeque, R. J. (2007). \emph{Finite Difference Methods for Ordinary
  and Partial Differential Equations}. SIAM.
  \texttt{https://doi.org/10.1137/1.9780898717839}. Convergence-order
  analysis for the periodic-trapezoidal-quadrature midpoint rule used
  by \passthrough{\lstinline!func::integrate!}.
\end{itemize}
